\scanpage{29}

d) \uline{Construction d'un schéma abélien formel}.

\diagram{\begin{tikzcd}[ampersand replacement=\&,column sep=small,row sep=small]
{} \& {} \& \widehat A' \& \widehat A'' \& \widehat A''' \\
M\arrow[d] \& S_0\arrow[l] \& S'_0\arrow[l] \& S''_0\arrow[l,shift left=.55ex]\arrow[l,shift right=.55ex] \& {} \\
T \& \widehat S\arrow[l] \& \widehat S'\arrow[l,dash] \& \widehat S''\arrow[l,shift left=.55ex]\arrow[l,shift right=.55ex] \& \widehat S'''\arrow[l,shift left=1ex]\arrow[l,dash]\arrow[l,shift right=1ex] \\
{} \& S_0 \& S'_0\arrow[l,dash] \& S''_0\arrow[l,dash,shift left=.5ex]\arrow[l,dash,shift right=.5ex] \& S'''_0\arrow[l,dash,shift left=1ex]\arrow[l,dash]\arrow[l,dash,shift right=1ex]
\end{tikzcd}}

Notons d'abord que si $S_0=\operatorname{Spec} k$ ($k$ corps résiduel) alors $\deleted{\unsure{B|A}}S'_0/S_0$ a une section, donc $S'_0\unsure{-|\to}S_0$ est un morphisme à descente effective universelle, et on trouve un $A_0$ sur $\deleted{\unsure{A}}\unsure{S|\Delta}_0$ qui descend $A'_0$.

Commençons par \deleted{prendre} les complétés formels de $S',S'',S''',A',A'',A'''$ \deleted{\unsure{sur}} relativement aux fibres de $s$, on trouve \deleted{\unsure{un|des}} schéma abélien formel $\widehat A'$ sur $\widehat S'$, avec donnée de descente pour $\widehat S'-\widehat S$. On va commencer \unsure{à|:} descendre dans la catégorie des schémas abéliens formels.

On introduit $T=\operatorname{Spf} W(k)$ ($W$ vecteurs de Witt) et $M={}$schéma formel des modules de $A_0$ ($M$ considéré comme schéma formel). Un lemme facile, qui est pratiquement trivial si $S'$ fini sur $S$ [donc $\widehat S',\widehat S''$, \unsure{\&c.|etc.} sont pro-artiniens] dit que la donnée d'une extension \inserted{formelle} $\widehat A'$ de $A'_0=A_{0\,S'_0}$ [resp. $\widehat A''_0$ de $A''_0$.] \&c. équivaut à la donnée d'un $T$-morphisme $\varphi':\widehat S'\to M$ [resp. $\varphi'':\widehat S''\to M$, \deleted{resp. $\widehat S'''-M$}] \&c. La donnée de descente \inserted{sur $\widehat A'$} \unsure{prolongeant} celle que l'on avait sur $A'_0$ équivaut à la condition $\varphi'\widehat p_1=\varphi'\widehat p_2$. Or on constate aussitôt (\unsure{relemme}, trivial encore en \unsure{le|si} $S'$ fini sur $S$, et \unsure{provenant|recourant} du th. de comparaison formel-algébrique en général) que la suite $\widehat S\longleftarrow\widehat S'\mathrel{\substack{\longleftarrow\\[-.6ex]\longleftarrow}}\widehat S''\deleted{\unsure{\leftleftarrows|\Leftarrow}}$ des schémas formels est exacte, donc on trouve $\varphi:\widehat S-M$ t.q. que $\varphi\widehat f=\varphi'$. Cela signifie qu'on a par descente
% archive sheet number: 28
